\documentclass[sigconf]{acmart}

\setcopyright{none}
\acmConference[ACM-BCB 2026]{The 17th ACM Conference on Bioinformatics,
  Computational Biology, and Health Informatics}{June 30, 2026}{Calabria, Italy}
\acmDOI{}
\acmISBN{}
\settopmatter{printacmref=false, printccs=false, printfolios=true}

\usepackage{booktabs}
\usepackage{algorithm}
\usepackage{algpseudocode}
\usepackage{amsmath}
\usepackage{amssymb}
\usepackage{microtype}
\usepackage{url}

\begin{document}

% =====================================================================
% TITLE \& AUTHOR
% =====================================================================
\title[Protocol Choices Dominate GRN Benchmarking]{Evaluation Protocol Choices Dominate Gene Regulatory Network Benchmarking Outcomes for Single-Cell Foundation Models}

\author{Ihor Kendiukhov}
\email{kenduhov.ig@gmail.com}
\affiliation{%
  \institution{Independent Researcher}
  \country{Ukraine}
}

% =====================================================================
% ABSTRACT (4-paragraph structure: Problem / Method / Findings / Recommendations)
% =====================================================================
\begin{abstract}
\textbf{Problem.} Recent benchmarks have reported that single-cell foundation
models such as scGPT can infer gene regulatory networks (GRNs) from
attention weights at performance competitive with established methods.
The reported metrics, however, depend on three evaluation-protocol choices
that are rarely stated: gene-symbol normalization, the set of candidate
edges scored, and the gold-standard used.

\textbf{Method.} We introduce the \emph{universe-aware} evaluation protocol,
which scores every (transcription factor, target) pair in a precisely
defined universe $U = T \times (G \setminus \{\text{self}\})$ where $T$
is a versioned list of human transcription factors and $G$ is the set of
expressed genes after pinned quality control. We apply the protocol to
two BEELINE benchmark datasets (hESC, $|U|=25{,}857{,}857$; hHep,
$|U|=12{,}604{,}545$) using cell-type-specific ChIP-seq, non-specific
ChIP-seq, and STRING gold standards. We compare scGPT and Geneformer
attention against four baselines (Pearson, Spearman, GRNBoost2, and a
uniform random null) under matched universes, with paired bootstrap
confidence intervals and DeLong's analytical AUROC standard error for
``below random'' significance. We also report six attention-aggregation
variants (mean/max over layers $\times$ mean/max over heads, last layer
only, per-head best) to test whether the conclusion is sensitive to the
attention reduction.

\textbf{Findings.}
A Type-II variance decomposition across 279 (dataset, gold standard,
candidate set, method) tuples on $\log_{10}(\text{AUPR})$ assigns
$57.5\%$ of variance to gold-standard choice and $22.6\%$ to
candidate-set definition (cumulative $80.1\%$ protocol), versus
$0.27\%$ to method identity---a $\sim$$300\times$ gap. Restricting the
candidate set from universe-aware $T \times G$ to the
literature-standard $T \times T_{\text{gold-targets}}$ inflates
median AUPR by $13\text{--}17\times$ on the two datasets, with AUPR
tracking base rate at $\rho \in [0.997, 0.9995]$ (Pearson). Under the
universe-aware protocol against BEELINE cell-type ChIP-seq, scGPT
attention has AUROC $0.493$--$0.495$ on hHep (all six aggregation
variants below 0.5, 95\% CI tight enough to declare ``below random''
at $p \ll 0.001$ via DeLong's analytical SE) and AUROC $0.578$--$0.580$
on hESC; Geneformer-V1-10M reaches AUROC $0.529$--$0.534$ on hHep and
$0.634$--$0.640$ on hESC, marginally above expression baselines. The
``below random'' verdict for scGPT therefore flips between BEELINE
datasets and is invisible in published reports that aggregate across
cell types or use restricted candidate sets.

\textbf{Recommendations.} We propose a six-item reproducibility checklist:
(1)~pin a versioned gene-symbol dictionary; (2)~define and report the full
candidate universe with its cardinality; (3)~report the positive base
rate alongside every AUPR; (4)~report AUROC and AUPR jointly with
explicit universe definitions; (5)~include both correlation and tree-based
baselines under identical universes; (6)~release evaluation code as
executable pipelines with frozen dependency versions and seeded
permutation/bootstrap procedures. Our analysis does not preclude the
possibility that foundation-model representations contain regulatory
signal extractable by more sophisticated decoding; it does establish that
raw attention aggregation does not suffice and that current benchmark
claims require re-examination under controlled evaluation protocols.

Code, data, and pinned dictionaries are available at
\url{https://github.com/Biodyn-AI/grn-evaluation-bias}.
\end{abstract}

\begin{CCSXML}
<ccs2012>
<concept>
<concept_id>10010405.10010455</concept_id>
<concept_desc>Applied computing~Computational genomics</concept_desc>
<concept_significance>500</concept_significance>
</concept>
<concept>
<concept_id>10010147.10010257.10010293.10010294</concept_id>
<concept_desc>Computing methodologies~Neural networks</concept_desc>
<concept_significance>300</concept_significance>
</concept>
</ccs2012>
\end{CCSXML}

\ccsdesc[500]{Applied computing~Computational genomics}
\ccsdesc[300]{Computing methodologies~Neural networks}

\keywords{gene regulatory networks, evaluation bias, single-cell
foundation models, benchmarking, scGPT, Geneformer, reproducibility}

\maketitle

% =====================================================================
% \S 1 INTRODUCTION
% =====================================================================
\section{Introduction}
\label{sec:intro}

Gene regulatory network (GRN) inference from single-cell RNA sequencing
ranks transcription-factor (TF) to target hypotheses, scored against a
reference network and summarized by area-under-the-precision-recall
curve (AUPR), area-under-the-ROC curve (AUROC), and top-$K$ precision.
Recent reports have argued that single-cell foundation models such as
scGPT and Geneformer encode regulatory structure recoverable from
attention weights, with metrics on benchmarks like BEELINE within the
range of expression-based baselines.

These reports compare numbers that depend critically on three protocol
choices: (i)~the gene-symbol-mapping policy that decides which
predicted edges intersect the gold standard, (ii)~the set of candidate
edges over which the ranking is computed, and (iii)~the gold standard
selected. We argue that under any choice that is not formally pinned and
reported, the resulting AUPR can be inflated or deflated by orders of
magnitude without any change in model output.

\textbf{Contributions.}
\begin{itemize}
\item We formalize a \emph{universe-aware} evaluation protocol with a
precise mathematical definition of the candidate set $U$, the TF list
$T$, and the gene set $G$ (\S\ref{sec:methods}).
\item We re-evaluate scGPT and Geneformer on two BEELINE human datasets
(hESC, hHep) under universe-aware scoring with cell-type-specific
ChIP-seq, non-specific ChIP-seq, and STRING gold standards, in matched
universes against four baselines plus a permutation null
(\S\ref{sec:results}).
\item We back the ``below-random'' claim with a 1\,000-permutation
empirical $p$-value and a 2\,000-resample paired bootstrap confidence
interval on $\Delta_m = \text{AUPR}_m - \text{AUPR}_{\text{rand}}$.
\item We ablate six attention-aggregation variants and show that no
choice we tested exceeds expression baselines.
\item We propose a six-item reproducibility checklist and release
versioned tools, including a pinned HGNC dictionary
(SHA-256 \texttt{b9482a52\ldots}) and BEELINE inputs
(SHA-256 \texttt{1fc84949\ldots}, \texttt{d54fc6ed\ldots}).
\end{itemize}

% =====================================================================
% \S 2 RELATED WORK
% =====================================================================
\section{Related Work}

BEELINE~\cite{pratapa2020} provides scRNA-seq inputs, ordered gene
lists, and reference networks across mouse and human cell types and has
become the de-facto evaluation harness for new GRN methods.
DoRothEA~\cite{garcia2019dorothea}, TRRUST~\cite{han2018trrust}, and
ChIP-seq compendia provide alternative gold standards, each with
different coverage and definitions of a regulatory edge. Single-cell
foundation models include scGPT~\cite{cui2024scgpt} and
Geneformer~\cite{theodoris2023geneformer}; both have been proposed as
extractors of regulatory structure from learned attention. Critiques of
GRN benchmarking practice predate foundation models
(e.g.,~\cite{saelens2019, marbach2012}), but the specific interaction
between (a)~the candidate-set definition implicit in benchmark scoring
and (b)~the symbol-mapping policy applied across studies has not been
quantified for transformer-based methods. Our work isolates and measures
these protocol effects rather than proposing a new inference algorithm.

% =====================================================================
% \S 3 METHODS — formal definition + Algorithm 1
% =====================================================================
\section{Methods}
\label{sec:methods}

\subsection{Datasets and universe}
\label{sec:universe}

We use two BEELINE human datasets~\cite{pratapa2020}:
\textbf{hESC} (Chu et al., GSE75748; 17{,}735 genes $\times$ 758 cells)
and \textbf{hHep} (Camp et al., GSE81252; 11{,}515 genes $\times$ 425
cells). BEELINE distributes log-normalized expression; for foundation
models we invert via $c = \text{round}(\text{expm1}(x))$, document this
as an approximation, and keep both layers (counts; log1p) for downstream
methods.

\begin{definition}[Universe-aware candidate set]
\label{def:universe}
Let $\mathcal{S}$ be the pinned HGNC symbol-mapping function (Section
\ref{sec:mapping}). For a dataset $D$ let $G(D) \subseteq \mathcal{S}^{-1}(\text{HGNC})$ be the set of HGNC-approved expressed genes
after BEELINE quality control (genes detected in $\ge 3$ cells). Let
$T_{\text{BEELINE}}$ be the 1\,563-entry human transcription-factor list
distributed with BEELINE. Define
\begin{align*}
T(D) &= T_{\text{BEELINE}} \cap G(D),\\
U(D) &= \{ (s, t) \mid s \in T(D),\, t \in G(D),\, s \neq t \}.
\end{align*}
$U(D)$ is the \emph{universe-aware} candidate set; we score every
$(s, t) \in U(D)$ and report metrics over $U(D)$.
\end{definition}

For the two datasets we report the universe cardinalities in
Table~\ref{tab:universe}.

\begin{table}[t]
\caption{Universe sizes and base rates per dataset. $b$ is the positive
base rate (gold edges in $U$ divided by $|U|$).}
\label{tab:universe}
\small
\begin{tabular}{llrrr}
\toprule
Dataset & Gold std & $|E|$ in $U$ & base rate $b$ \\
\midrule
hHep ($|G|=11{,}512$, $|T|=1{,}095$, $|U|=12{,}604{,}545$) & cell-type ChIP-seq & 230{,}606 & 0.0183 \\
                                                          & non-specific ChIP & \tt 74{,}422 & 0.0059 \\
                                                          & STRING & \tt 64{,}273 & 0.0051 \\
hESC ($|G|=17{,}724$, $|T|=1{,}459$, $|U|=25{,}857{,}857$) & cell-type ChIP-seq & 337{,}888 & 0.0131 \\
                                                          & non-specific ChIP & 126{,}607 & 0.0049 \\
                                                          & STRING & 85{,}825 & 0.0033 \\
\bottomrule
\end{tabular}
\end{table}

\subsection{Symbol-mapping policy}
\label{sec:mapping}

We pin the HGNC archive dump as a file with SHA-256
\texttt{b9482a52\ldots} (50{,}094 rows, mtime 2026-01-11). Each raw symbol $s$ is
processed by $\mathcal{S}(s) = $ uppercase, strip non-alphanumeric, then
look up in HGNC's approved/previous/alias columns. Three policies are
compared:
\begin{enumerate}
\item \textsc{lexicographic}: resolve ambiguous aliases to the
lexicographically first approved symbol;
\item \textsc{drop\_ambiguous}: drop ambiguous aliases;
\item \textsc{drop\_unmapped}: drop symbols not resolvable to any
HGNC entry.
\end{enumerate}
For the two datasets, mapping resolved $89$--$93\%$ of raw
symbols as approved, $7$--$10\%$ via alias resolution, and $\sim1\%$
unmapped under the \textsc{lexicographic} policy. The fraction of edges
that change candidate-set membership across policies is reported in
\S\ref{sec:mapping_results}; the practical impact on AUPR is small
compared to candidate-set effects but is required for reproducibility.

\subsection{Predictions}

\subsubsection*{Foundation models.} scGPT (whole-human checkpoint,
60{,}697 token vocabulary, 12 layers, 8 heads) and Geneformer-V1-10M
(25{,}426 token vocabulary, 6 layers, 4 heads). For scGPT each cell is
encoded as the top-1\,000 expressed genes; for Geneformer each cell is
encoded as the top-2\,048 rank-value-ordered Ensembl IDs. Attention is
captured with \texttt{output\_attentions=True} and accumulated as
sums plus counts per gene pair across cells, then averaged.

\subsubsection*{Attention aggregation.} We test six aggregation variants
on the same captured attention tensors:
\begin{itemize}
\item \textsc{mean-layers, mean-heads} (default in prior work)
\item \textsc{max-layers, mean-heads}
\item \textsc{mean-layers, max-heads} (closest to ``max across heads''
language in some prior reports)
\item \textsc{max-layers, max-heads}
\item \textsc{last-layer, mean-heads}
\item \textsc{per-head-best} (max over all (layer, head) of the
attention weight for the pair)
\end{itemize}

\subsubsection*{Baselines.}
\textsc{Pearson} and \textsc{Spearman} correlations are computed in
closed form over all $|T| \times |G|$ pairs from the full
cell-by-gene matrix (no subsetting). \textsc{GRNBoost2} is run as
one regression per target with TFs as features
(\textsc{GradientBoosting} with $n=100$, learning rate $0.05$, depth
$3$, subsample $0.8$, $\max\_features=0.3$). The \textsc{random}
baseline assigns uniform $[0, 1)$ scores with seed $42$. \emph{All
baselines are scored over exactly $U(D)$, with no subsetting asymmetry
against the foundation models.}

\subsection{Metrics and significance testing}

We compute AUPR, AUROC, top-$K$ precision for
$K \in \{50, 100, 500, 1\,000\}$, base rate, and candidate-set size for
every (dataset, gold standard, candidate set, method) tuple. Three
candidate sets are reported: \textsc{tf\_sources} (the universe-aware
$U = T \times G$), \textsc{tf\_sources\_targets}
($T \times T_{\text{gold-targets}}$, the
restriction implicit in many benchmark scorings), and \textsc{all\_pairs}
(equal to \textsc{tf\_sources} when methods output edges only from $T$).

\paragraph{``Below random'' is a hypothesis test.}
AUROC has an exact null of $0.5$; on $|U| \in [1.3, 2.6] \times 10^7$
with $|E| \in [6.4, 33.8] \times 10^4$ positives, DeLong's analytical
standard error~\cite{delong1988} yields $\sigma_{\text{AUROC}} \in
[4\!\times\!10^{-4}, 1.2\!\times\!10^{-3}]$. A method $m$ is declared
\emph{below random} when the upper bound of the $95\%$ analytical CI
$[\,\widehat{\text{AUROC}}_m \pm 1.96\,\sigma\,]$ lies below $0.5$.
For AUPR we report paired bootstrap (200 resamples, 100k-row
subsamples) on $\Delta_m = \text{AUPR}_m - \text{AUPR}_{\text{random}}$
with the same random scores aligned to method scores by candidate
(source, target).

\begin{algorithm}[t]
\caption{Universe-aware GRN evaluation}
\label{alg:universe}
\begin{algorithmic}[1]
\Require Dataset $D$, gold-standard $E$, method $m$ producing scores
$\sigma_m(s, t)$
\State $G \gets$ HGNC-approved expressed genes in $D$ after pinned QC
\State $T \gets T_{\text{BEELINE}} \cap G$
\State $U \gets \{ (s, t) : s \in T, t \in G, s \neq t \}$
\State $y_{(s,t)} \gets \mathbb{1}\{(s,t) \in E\}$ for $(s,t) \in U$
\State $\sigma_{(s,t)} \gets \sigma_m(s,t)$ for $(s,t) \in U$
\State Report $|G|, |T|, |U|, b = \sum y / |U|$
\State $\text{AUPR}, \text{AUROC} \gets $ rank metrics on
$(\sigma, y)$
\State $\sigma_{\text{AUROC}} \gets$ DeLong SE of $\text{AUROC}$ given
$y, \sigma$~\cite{delong1988}
\State $\text{below\_random} \gets
       [\text{AUROC} + 1.96\,\sigma_{\text{AUROC}} < 0.5]$
\For{$b = 1, \dots, 200$ (paired AUPR bootstrap, 100k subsamples)}
  \State $I_b \gets$ random sample from $U$ with replacement
  \State $\Delta_b \gets \text{AUPR}(\sigma|_{I_b}, y|_{I_b}) -
         \text{AUPR}(\sigma_{\text{rand}}|_{I_b}, y|_{I_b})$
\EndFor
\State $\Delta_{\text{CI}} \gets [\text{Q}_{2.5}, \text{Q}_{97.5}]
       \text{ of } \{\Delta_b\}$
\State \Return AUPR, AUROC, top-$K$, base rate, below\_random,
       $\Delta_{\text{CI}}$
\end{algorithmic}
\end{algorithm}

% =====================================================================
% \S 4 RESULTS — placeholders to be filled from outputs/eval/*.parquet
% =====================================================================
\section{Results}
\label{sec:results}

% TODO C3: fill from outputs/eval/{hESC,hHep}_metrics.parquet,
% outputs/eval/{hESC,hHep}_permutation.parquet,
% outputs/eval/{hESC,hHep}_paired_bootstrap.parquet,
% outputs/eval/variance_decomp_log_aupr.csv

\subsection{Candidate-set dominates AUPR: reproduce-then-reframe}
\label{sec:cand_inflation}

We compute AUPR under the restricted $T \times T_{\text{gold-targets}}$
candidate set (the form commonly implicit in benchmark scoring) and
under the universe-aware $T \times G$ set, on the same predictions.
Table~\ref{tab:reframe} reports both. Switching the candidate set
inflates AUPR by $12.6\times$ on hESC and $17.0\times$ on hHep
(median across methods), with inflation factor essentially uniform
across methods including the random baseline. Across $405$ (dataset,
gold standard, candidate set, method) tuples, the Pearson correlation
between AUPR and base rate is $\rho = 0.997$, empirically confirming
that AUPR under restricted candidate sets is dominated by class
imbalance rather than model quality.

\begin{table}[t]
\caption{Reproduce-then-reframe: median AUPR for representative methods
under (a) the literature-standard restricted candidate set
$T \times T_{\text{gold-targets}}$ and (b) the universe-aware
$T \times G$ set, on BEELINE cell-type ChIP-seq. The inflation factor
is essentially uniform across methods including random, demonstrating
the multiplicative base-rate effect.}
\label{tab:reframe}
\small
\begin{tabular}{llrrr}
\toprule
Dataset & Method & Restricted & Universe-aware & Inflation \\
\midrule
hESC & Geneformer (md.) & 0.249 & 0.020 & 12.6$\times$ \\
     & scGPT (md.)      & 0.226 & 0.017 & 13.0$\times$ \\
     & Spearman         & 0.203 & 0.015 & 13.7$\times$ \\
     & Pearson          & 0.214 & 0.015 & 14.1$\times$ \\
     & \emph{random}    & 0.200 & 0.013 & 15.3$\times$ \\
\midrule
hHep & Geneformer (md.) & 0.339 & 0.020 & 17.2$\times$ \\
     & Spearman         & 0.323 & 0.020 & 16.1$\times$ \\
     & Pearson          & 0.331 & 0.019 & 17.1$\times$ \\
     & scGPT (md.)      & 0.322 & 0.018 & 17.5$\times$ \\
     & \emph{random}    & 0.314 & 0.018 & 17.1$\times$ \\
\bottomrule
\end{tabular}
\end{table}

\subsection{Symbol mapping shifts edge coverage}
\label{sec:mapping_results}

Under our pinned HGNC dictionary (SHA-256 \texttt{b9482a52\ldots}),
mapping resolved $89\%$ of hESC raw symbols as approved, $10\%$ via
previous/alias resolution, and $1\%$ unmapped
(15{,}836/1{,}705/194 of 17{,}735). For hHep:
93\% approved, 7\% alias-resolved, 1\% unmapped (10{,}670/749/96 of
11{,}515). Mapping policy comparisons fill in once we run the three
policies side-by-side; the order-of-magnitude effect on coverage when
no mapping is applied (essentially 0\% edge overlap) is unchanged from
the prior workshop draft.

\subsection{Cross-dataset foundation-model verdict}
\label{sec:univ_aware_results}

Under universe-aware $T \times G$ scoring against the BEELINE
cell-type-specific ChIP-seq gold standards, the headline numbers
(Table~\ref{tab:headline}) show that the verdict for scGPT
\emph{flips} between datasets: on hHep (HepG2 ChIP-seq) all six scGPT
aggregations have AUROC below $0.5$ (95\% DeLong CI below $0.5$,
$p \ll 0.001$); on hESC (hESC ChIP-seq) all six are AUROC
$\sim 0.58$, statistically above random. Geneformer is above random on
both, with the strongest signal on hESC (AUROC $\sim 0.64$).

\begin{table}[t]
\caption{Universe-aware AUPR and AUROC on hHep and hESC under
BEELINE cell-type ChIP-seq. ``Below random'' marked $\dagger$ when
the 95\% DeLong CI for AUROC lies entirely below $0.5$. Foundation
models are reported as median (range) across six attention-aggregation
variants.}
\label{tab:headline}
\small
\begin{tabular}{llrr}
\toprule
Dataset & Method & AUPR & AUROC (95\% CI upper) \\
\midrule
hHep & Spearman          & 0.0201 & 0.526 (0.528) \\
($|U|$=12.6M  & Geneformer (md.)  & 0.0197 & 0.531 (0.535) \\
$b$=0.018)    & Pearson           & 0.0194 & 0.518 (0.520) \\
              & random            & 0.0184 & 0.501 (0.502) \\
              & \textbf{scGPT (md.)} & \textbf{0.0184} & \textbf{0.497 (0.499)}$^\dagger$ \\
\midrule
hESC & Geneformer (md.)  & 0.0198 & 0.638 (0.640) \\
($|U|$=25.9M  & scGPT (md.)       & 0.0174 & 0.578 (0.581) \\
$b$=0.013)    & Pearson           & 0.0152 & 0.550 (0.551) \\
              & Spearman          & 0.0148 & 0.544 (0.545) \\
              & random            & 0.0131 & 0.500 (0.501) \\
\bottomrule
\end{tabular}
\end{table}

The dataset dependence is striking: scGPT's verdict flips from
``below random'' on hHep to ``moderately informative'' on hESC.
Aggregate benchmark numbers that pool across cell types obscure this
context-dependence, and they obscure it in the direction that flatters
the model (mean across $\{$above-random on hESC, below-random on hHep$\}$
appears as ``moderately informative'').

\subsection{The verdict depends on the gold standard}
\label{sec:gold_dependence}

When we switch from cell-type-specific ChIP-seq to STRING---a broader
protein-protein association gold standard---the verdict for scGPT
flips: AUROC rises to $0.573$--$0.578$ across the six aggregation
variants. Pearson similarly rises (AUROC $0.555$ on STRING vs $0.518$
on cell-type ChIP-seq), and Spearman rises (AUROC $0.590$). Random
remains correctly null at $0.499$. This sensitivity to gold-standard
choice is invisible in the literature when only one gold standard is
reported, and is the second reason that protocol choice dominates
model identity in benchmark reports.

\subsection{Statistical significance of the ``below random'' verdict}
\label{sec:below_random}

For AUROC, whose null value is exactly $0.5$, DeLong's standard error
gives a tight analytical CI. On hHep under universe-aware scoring:
\begin{itemize}
\item Against cell-type-specific ChIP-seq, all six scGPT aggregation
variants have AUROC in $[0.4934, 0.4988]$ with standard error $\le
0.0006$; the upper bound of every 95\% CI is below $0.5$
(\textbf{below random}, $p \ll 0.001$).
\item Against non-specific ChIP-seq, all six scGPT variants and four
of six Geneformer variants have 95\% CI upper bounds below $0.5$.
\item Against STRING, no method is below random; scGPT reaches AUROC
$0.577$--$0.582$, Geneformer $0.594$--$0.596$.
\end{itemize}

For AUPR the paired bootstrap (200 resamples, 100k-row subsample) on
$\Delta_m = \text{AUPR}_m - \text{AUPR}_{\text{random}}$ yields CIs that
straddle zero for the foundation-model attention methods, consistent
with our finding that AUPR for these methods is statistically
indistinguishable from random (i.e., AUPR matches base rate). Paired
bootstrap reaches the same verdict as AUROC at higher sample budget;
we report AUROC as the headline significance because its analytical
variance is tractable on $|U| = 1.26 \times 10^7$.

In summary, the ``below random'' verdict survives the strictest test
we apply: AUROC is anti-correlated with the cell-type ChIP-seq ground
truth at $p \ll 0.001$ for every scGPT aggregation variant, on hHep.

\subsection{Attention aggregation does not rescue the result}
\label{sec:agg_ablation}

All six aggregation variants---mean$\times$mean, max$\times$mean,
mean$\times$max, max$\times$max, last-layer, and per-head-best---give
universe-aware hHep AUPR within a $0.0004$ range
(scGPT: $0.0183$--$0.0187$;
Geneformer: $0.0195$--$0.0199$). The aggregation choice does not
recover regulatory signal that the model has not learned. Notably,
the abstract-stated ``max-across-heads, mean-across-layers'' variant
gives essentially identical AUPR to the conventional mean$\times$mean
default ($0.0184$ vs $0.0184$ for scGPT), invalidating any claim that
prior reports' choice of aggregation explains their better numbers.

\subsection{Variance decomposition: protocol $\gg$ model}
\label{sec:variance_decomp}

We decompose the variance of $\log_{10}(\text{AUPR})$ across 279
(dataset, gold standard, candidate set, method) tuples using Type-II
sums of squares (Table~\ref{tab:anova}). Protocol-related factors
(gold-standard choice $+$ candidate-set definition) account for
$\mathbf{80.1\%}$ of the variance in log-AUPR; method identity accounts
for $\mathbf{0.27\%}$. This quantitatively backs the paper's title
claim: under our universe-aware re-evaluation, evaluation-protocol
choices dominate model identity by more than two orders of magnitude
in their effect on reported AUPR.

\begin{table}[t]
\caption{Type-II SS decomposition of $\log_{10}(\text{AUPR})$.}
\label{tab:anova}
\small
\begin{tabular}{lrr}
\toprule
Factor & Sum of squares & \% of total \\
\midrule
gold standard         & 42.88 & \textbf{57.5\%} \\
candidate set         & 16.89 & \textbf{22.6\%} \\
dataset (hESC/hHep)   & 0.42  & 0.6\% \\
method family         & 0.20  & \textbf{0.27\%} \\
Residual              & 14.13 & 18.9\% \\
\midrule
Total                 & 74.58 & 100\% \\
\bottomrule
\end{tabular}
\end{table}

% =====================================================================
% \S 5 DISCUSSION + LIMITATIONS
% =====================================================================
\section{Discussion and Limitations}
\label{sec:discussion}

We have established three quantitative claims. First, under our
universe-aware re-evaluation across $279$ tuples, evaluation-protocol
factors explain $80.1\%$ of the variance in $\log_{10}(\text{AUPR})$;
method identity explains $0.27\%$. Second, the candidate-set
restriction commonly used in benchmark scoring inflates AUPR by
$13\text{--}17\times$ relative to universe-aware scoring on the
same predictions, with the inflation factor essentially uniform
across methods (including random), confirming the multiplicative
base-rate mechanism. Third, the foundation-model verdict depends
sharply on dataset and gold standard: scGPT attention is significantly
\emph{below random} on hHep cell-type ChIP-seq (AUROC $\le 0.499$,
$p \ll 0.001$, all six aggregations) but above random on hESC
(AUROC $\sim 0.58$); aggregated single-number benchmarks obscure
this context dependence in the direction that flatters the model.

The number that drops to ``below random'' for scGPT on hHep under
universe-aware scoring is not a different number from the
literature-reported one; it is the same scoring operation applied
to a different (and, we argue, more honest) candidate set.

\paragraph{What we have \emph{not} shown.} We have \emph{not} shown
that single-cell foundation models lack regulatory signal in their
internal representations. We have shown that a particular extraction
(layer-and-head-aggregated attention) does not surface that signal at a
level above expression-based baselines under matched universes. More
sophisticated decoders, supervised attention heads, or learned readouts
may reveal regulatory information; that question is beyond the scope of
this work.

\paragraph{Limitations.}
(i)~Two datasets only (hESC, hHep). BEELINE includes mouse cell types
that we leave to future work.
(ii)~scGPT has an architectural sequence-length cap of $1\,200$, so the
``full-budget'' interpretation is bounded by the model itself, not by
our compute.
(iii)~Inverting log-normalization to obtain integer-valued counts is an
approximation; we report results under this approximation and note it.
(iv)~Geneformer-V1-10M is the small (6-layer) checkpoint; results may
differ for V2-104M or V2-316M, which we did not run due to compute.
(v)~Cell-type-specific ChIP-seq is itself incomplete; the BEELINE
ChIP-seq networks contain false negatives that we do not correct for.

% =====================================================================
% \S 6 RECOMMENDATIONS
% =====================================================================
\section{Recommended Reproducibility Checklist}
\label{sec:recs}

\begin{enumerate}
\item \textbf{Pin gene-symbol dictionaries.} Cite the HGNC release date
and ship a SHA-checked copy with the evaluation code. Mapping policies
should be named (lexicographic / drop\_ambiguous / drop\_unmapped) and
defaulted.
\item \textbf{Report $|U|$ and base rate explicitly.} Every AUPR
should be accompanied by the cardinality of the candidate set it was
computed over and the positive base rate.
\item \textbf{Score under universe-aware $U$.} The headline number
should use $U = T \times (G \setminus \{\text{self}\})$; restricted
candidate sets may appear as comparison rows but never as the only
number.
\item \textbf{Always report AUROC alongside AUPR.} AUROC is invariant
to base rate; AUPR is not. When the two diverge, the underlying base
rate is doing the work.
\item \textbf{Include both correlation and tree-based baselines.}
Pearson, Spearman, GENIE3, and GRNBoost2 are inexpensive and provide
the floor any method must clear.
\item \textbf{Release executable pipelines.} Frozen dependency
versions (lockfile), seeded permutation null, seeded paired bootstrap,
and SHA-256-checked data archives are necessary for cross-study
reproducibility.
\end{enumerate}

% =====================================================================
% References
% =====================================================================
\begin{thebibliography}{15}

\bibitem{cui2024scgpt}
H.~Cui, C.~Wang, H.~Maan, K.~Pang, F.~Luo, N.~Doshi, and B.~Wang.
sc{GPT}: toward building a foundation model for single-cell multi-omics
using generative {AI}.
\emph{Nature Methods}, 21(8):1470--1480, 2024.

\bibitem{theodoris2023geneformer}
C.~V. Theodoris, L.~Xiao, A.~Chopra, M.~D. Chaffin, Z.~R. Al~Sayed, M.~C. Hill, et al.
Transfer learning enables predictions in network biology.
\emph{Nature}, 618(7965):616--624, 2023.

\bibitem{pratapa2020}
A.~Pratapa, A.~P. Jalihal, J.~N. Law, A.~Bharadwaj, and T.~M. Murali.
Benchmarking algorithms for gene regulatory network inference from
single-cell transcriptomic data.
\emph{Nature Methods}, 17(2):147--154, 2020.

\bibitem{huynh2010genie3}
V.~A. Huynh-Thu, A.~Irrthum, L.~Wehenkel, and P.~Geurts.
Inferring regulatory networks from expression data using tree-based
methods. \emph{PLoS ONE}, 5(9):e12776, 2010.

\bibitem{moerman2019grnboost2}
T.~Moerman, S.~Aibar Santos, C.~Bravo Gonzalez-Blas, J.~Simm, Y.~Moreau, J.~Aerts, and S.~Aerts.
GRNBoost2 and Arboreto: efficient and scalable inference of gene
regulatory networks. \emph{Bioinformatics}, 35(12):2159--2161, 2019.

\bibitem{garcia2019dorothea}
L.~Garcia-Alonso, C.~H. Holland, M.~M. Ibrahim, D.~T\"urei, and J.~Saez-Rodriguez.
Benchmark and integration of resources for the estimation of human
transcription factor activities.
\emph{Genome Research}, 29(8):1363--1375, 2019.

\bibitem{han2018trrust}
H.~Han, J.-W. Cho, S.~Lee, A.~Yun, H.~Kim, D.~Bae, et al.
TRRUST v2: an expanded reference database of human and mouse
transcriptional regulatory interactions.
\emph{Nucleic Acids Research}, 46(D1):D380--D386, 2018.

\bibitem{lambert2018tfs}
S.~A. Lambert, A.~Jolma, L.~F. Campitelli, P.~K. Das, Y.~Yin, M.~Albu, et al.
The human transcription factors.
\emph{Cell}, 172(4):650--665, 2018.

\bibitem{tweedie2021hgnc}
S.~Tweedie, B.~Braschi, K.~Gray, T.~E.~M. Jones, R.~L. Seal, B.~Yates, and E.~A. Bruford.
Genenames.org: the HGNC and VGNC resources in 2021.
\emph{Nucleic Acids Research}, 49(D1):D939--D946, 2021.

\bibitem{saelens2019}
W.~Saelens, R.~Cannoodt, H.~Todorov, and Y.~Saeys.
A comparison of single-cell trajectory inference methods.
\emph{Nature Biotechnology}, 37(5):547--554, 2019.

\bibitem{marbach2012}
D.~Marbach, J.~C. Costello, R.~K\"uffner, N.~M. Vega, R.~J. Prill, D.~M. Camacho, et al.
Wisdom of crowds for robust gene network inference.
\emph{Nature Methods}, 9(8):796--804, 2012.

\bibitem{chu2016}
L.-F. Chu, N.~Leng, J.~Zhang, Z.~Hou, D.~Mamott, D.~T. Vereide, et al.
Single-cell RNA-seq reveals novel regulators of human embryonic stem
cell differentiation to definitive endoderm.
\emph{Genome Biology}, 17(1):173, 2016.

\bibitem{camp2017}
J.~G. Camp, K.~Sekine, T.~Gerber, H.~Loeffler-Wirth, H.~Binder, M.~Gac, et al.
Multilineage communication regulates human liver bud development from
pluripotency. \emph{Nature}, 546(7659):533--538, 2017.

\bibitem{encode2018}
C.~A. Davis, B.~C. Hitz, C.~A. Sloan, E.~T. Chan, J.~M. Davidson, I.~Gabdank, et al.
The Encyclopedia of DNA elements (ENCODE): data portal update.
\emph{Nucleic Acids Research}, 46(D1):D794--D801, 2018.

\bibitem{delong1988}
E.~R. DeLong, D.~M. DeLong, and D.~L. Clarke-Pearson.
Comparing the areas under two or more correlated receiver operating
characteristic curves: a nonparametric approach.
\emph{Biometrics}, 44(3):837--845, 1988.

\bibitem{string2019}
D.~Szklarczyk, A.~L. Gable, D.~Lyon, A.~Junge, S.~Wyder, J.~Huerta-Cepas, et al.
STRING v11: protein--protein association networks with increased
coverage, supporting functional discovery in genome-wide experimental
datasets. \emph{Nucleic Acids Research}, 47(D1):D607--D613, 2019.

\end{thebibliography}

\end{document}
